\section*{Part VI: Industry-Specific Pricing Systems}

\subsection*{45. Consumer Packaged Goods and Physical Products}

The price on a consumer shelf is not the manufacturer's realized revenue. Retailer margins, distributor discounts, trade promotions, returns, spoilage, and logistics matter. Consumers may make quick, habitual choices based on availability, pack size, sensory experience, and brand recognition rather than elaborate utility calculations. A premium product must survive repeated full-price purchase, not only launch sampling.

For appliances or industrial machinery, ownership cost: including repairs, spare parts, downtime, and warranty: can matter more than sticker price. Use total ownership arguments where they are testable, while modeling the cost of warranties and service commitments. A high advertised price with negative contribution after returns is not successful premium pricing.

\subsection*{46. Software and Enterprise AI}

SaaS models commonly use seat-based, tiered, usage, flat, or hybrid pricing. Enterprise buyers also demand predictability, security, approvals, and support. A price unit that seems optimal to the vendor may suppress user adoption if employees perceive every action as expensive. A platform charging by AI token may reward frugality rather than task completion.

Analyze the cost per successful workflow, quality, cloud and inference costs, model volatility, error review, and account support. Where the customer needs predictable expense, offer limits, included usage, clear overages, or contract terms. Annual contracts can finance development but create performance obligations and customer trust duties across the term.

\subsection*{47. Agencies, Consultancies, and Implementation Services}

Hourly work, fixed projects, retainers, licenses, and limited performance fees allocate risk differently. A specialist can legitimately price the value of diagnosis rather than number of hours, but the buyer needs a credible way to evaluate it. A fixed project fee must define dependencies and change control; otherwise uncontrolled scope destroys its economics.

A consultancy can command a premium by reducing client uncertainty through explicit acceptance criteria, implementation discipline, a specialized methodology, and documented results. ``We bring world-class experts'' is a weaker proposition than a clear diagnosis, accountable work plan, and transparent failure analysis. Price clarity can itself reduce the cost of purchasing the service.

\subsection*{48. Marketplaces, Platforms, and 2-Sided Pricing}

A marketplace must decide whether to charge buyers, sellers, advertisers, or some combination. Subsidizing 1 side can increase participation and liquidity for the other. Commission, subscription, listing, and advertising models distribute the value differently and affect the temptation to transact off-platform.

Optimizing take rate without measuring both sides can shrink the network. Model transaction frequency, matching reliability, disputes, multi-homing, leakage, and the actual reason participants pay. A registration count is not a pricing asset if transactions cannot be completed.

\subsection*{49. Research, Licensing, and Deep-Technology Products}

Some products have high fixed development and regulatory costs but low marginal replication cost; others require expensive manufacturing and technical support. Licensing can spread fixed research costs across buyers, but royalties require auditability and contract enforcement. Hardware may be sold at a margin, through long-term service, or as part of an integrated system.

Avoid using ``we spent 5 years inventing this'' as the price argument. The buyer pays for a useful, defensible advantage or access to something unavailable elsewhere. Development cost determines what the venture needs to recover; it does not automatically determine the buyer's willingness to pay.
